
\item
\uline{Pour} @@M0@@, @@M1@@ \uline{et} @@M2@@ \uline{sont plats sur} @@M3@@.
\item \uline{La suite}
@@M4@@
\uline{est exacte et} @@M5@@ \uline{est un} @@M6@@\uline{-module plat sur} @@M7@@.
\end{enumerate}
\uline{Alors} @@M8@@ \uline{est représentable par un morphisme quasi-affine de présentation finie.}

En effet notons d'abord que 1)-4) sont stables par tout changement de base @@M9@@. Seule la stabilité de 3) et 4) n'est pas évidente. Mais d'après la formule de Künneth (EGA III 6.9.9), 1) et 3) impliquent que la formation des @@M10@@ et @@M11@@ commute aux changements de base @@M12@@, ce qui implique aussitôt la stabilité des conditions 3). De plus, la formation du conoyau de @@M13@@ commute donc aussi aux changements de base~; comme il est plat sur @@M14@@ (d'après 4)), @@M15@@ est plat sur @@M16@@ et commute aux changements de base. Par suite, @@M17@@ commute aux changements de base @@M18@@, et est donc égal à @@M19@@ universellement.

La stabilité des hypothèses implique deux choses:a. L'assertion à démontrer étant locale pour f.p.p.f., on peut par descente supposer que @@M20@@ possède une section @@M21@@, donc @@M22@@ possède la section @@M23@@. b. Pour établir l'assertion il suffit de considérer un @@M24@@-point de @@M25@@, c'est-à-dire d'après a) et 2) un faisceau inversible @@M26@@ sur @@M27@@ muni d'une rigidification @@M28@@ relativement à @@M29@@, et de montrer que le faisceau f.p.p.f. @@M30@@ image réciproque de ce point dans @@M31@@ est représentable par un @@M32@@-schéma quasi-affine, de présentation finie sur @@M33@@.

